% !TEX program = xelatex
\documentclass[10pt,a4paper]{ctexart}

\usepackage{amsmath,amssymb}
\usepackage{booktabs}
\usepackage{geometry}
\usepackage{enumitem}
\usepackage{longtable}
\usepackage{array}
\usepackage{xcolor}
\usepackage{xurl}
\usepackage{hyperref}

\geometry{left=1.8cm,right=1.8cm,top=2.2cm,bottom=2.2cm}
\hypersetup{colorlinks=true,linkcolor=blue,urlcolor=blue}

\newcolumntype{L}[1]{>{\raggedright\arraybackslash}p{#1}}
\urlstyle{tt}
\newcommand{\eng}[1]{\begingroup\footnotesize\path{#1}\endgroup}
\newcommand{\modref}[2]{\textit{#1~#2}}

\title{%
  \textbf{星闪 SLB 物理层}\\[0.3em]
  \Large 6.5 模块实现参数汇总\\[0.2em]
  \large CRC/Gold/QPSK \textbullet\ 功控测量 \textbullet\ 调制 \textbullet\ 加扰/SC-FDMA/ZC
}
\author{依据 6.5 分册技术文档补全；airMode\,=\,\texttt{slb}\\
\small 规则：本模块参数列全；依赖他模块时只写其\textbf{输出}
}
\date{2026 年 7 月 28 日}

\begin{document}
\maketitle
\tableofcontents
\newpage

\section{约定}
本模块输出 / 直接入参 / 引用他模块输出（不展开上游内部）。

%----------------------------------------------------------------------
\part{6.5.1--6.5.3}
%----------------------------------------------------------------------

\section{6.5.1 CRC}
\textbf{API}：\eng{slbCalculateCrc} / \eng{slbValidateCrc}

\begin{longtable}{L{3.5cm}L{11.7cm}}
\caption{6.5.1}\label{tab:651} \\
\toprule \textbf{角色} & \textbf{内容} \\ \midrule
\endfirsthead
\toprule \textbf{角色} & \textbf{内容} \\ \midrule
\endhead
\bottomrule\endfoot
输出 &
\eng{crcEncodedBits[]}（$B{=}A{+}L$）；
\eng{crcPass}；
\eng{crcParityLengthL} \\
直接入参 &
\eng{infoBits[]}/\eng{infoBitCount}；
\eng{crcPoly}（24A/24B/16/8/6）；
条件 \eng{phyIdMaskBits}（仅 $L{=}24$） \\
他模块 &
\modref{6.2.4}{物理层控制信息}/\modref{6.2.5}{物理层数据} 已组装比特；
下游 \modref{6.2.6}{信道编码} 消费编码输出 \\
\end{longtable}

\section{6.5.2 Gold}
\textbf{API}：\eng{slbGenerateGoldSequence}

\begin{longtable}{L{3.5cm}L{11.7cm}}
\caption{6.5.2}\label{tab:652} \\
\toprule \textbf{角色} & \textbf{内容} \\ \midrule
\endfirsthead
\toprule \textbf{角色} & \textbf{内容} \\ \midrule
\endhead
\bottomrule\endfoot
输出 & \eng{goldBits[]}；$x_1$ 预计算表（实现可选） \\
直接入参 & \eng{cInit}、\eng{pnLength}（调用方已算好 $c_{\mathrm{init}}$） \\
L0 & 1600 偏移强制；$x_1$ 初态固定 \\
消费者 & \modref{6.5.3}{伪随机QPSK}（322）、\modref{6.5.7}{比特加扰}、\modref{6.5.9.1}{ZC组跳} \\
\end{longtable}

\section{6.5.3 伪随机 QPSK}
\textbf{API}：\eng{slbCalculateRsCinit} / \eng{slbGenerateRsQpsk} / \eng{slbMapRsQpskToRe}

\begin{longtable}{L{3.5cm}L{11.7cm}}
\caption{6.5.3}\label{tab:653} \\
\toprule \textbf{角色} & \textbf{内容} \\ \midrule
\endfirsthead
\toprule \textbf{角色} & \textbf{内容} \\ \midrule
\endhead
\bottomrule\endfoot
输出 &
\eng{rsCinit}；
\eng{rsQpskSequence[161]}；
按 $k$ 映射（$k{=}80\to0$） \\
直接入参 &
$n,l$；L1 \eng{gNodeIdLow8}；
\modref{6.2.2.1}{波形和符号}\eng{nCp}；
\modref{6.2.2.2}{无线帧和超帧}\eng{nSymFrame}；
\modref{6.5.2}{Gold序列}\eng{goldBits[]}（$M_{\mathrm{PN}}{=}322$） \\
禁混 & 不得与 \modref{6.5.6.1}{QPSK调制} API 混用 \\
\end{longtable}

%----------------------------------------------------------------------
\part{6.5.4--6.5.5}
%----------------------------------------------------------------------

\section{6.5.4 功率控制}
\textbf{API}：\eng{slbCalculatePathloss} / \eng{slbCalculateTxPower} / \eng{slbApplyPmaxLimit}

\begin{longtable}{L{3.5cm}L{11.7cm}}
\caption{6.5.4}\label{tab:654} \\
\toprule \textbf{角色} & \textbf{内容} \\ \midrule
\endfirsthead
\toprule \textbf{角色} & \textbf{内容} \\ \midrule
\endhead
\bottomrule\endfoot
输出 &
\eng{pathlossDb}；
\eng{txPowerTotalDbm}；
\eng{txPowerPerScLin[]}；
限幅总功率 \\
直接入参 &
L1 \eng{gLinkReferenceSignalPowerDbm}、各业务 \eng{*P0NominalDbm}、\eng{pMaxDbm}；
调度 $M,K$、多业务线性功率 \\
他模块 & \modref{6.5.5.1}{RSRP}\eng{rsrpDbm} \\
\end{longtable}

\section{6.5.5 测量量}

\begin{longtable}{L{2.2cm}L{4.5cm}L{8.5cm}}
\caption{6.5.5}\label{tab:655} \\
\toprule \textbf{子节} & \textbf{输出} & \textbf{入参/他模块} \\ \midrule
\endfirsthead
\toprule \textbf{子节} & \textbf{输出} & \textbf{入参/他模块} \\ \midrule
\endhead
\bottomrule\endfoot
RSRP & \eng{rsrpDbm} & 训练 RE 功率；\modref{6.2.3.1}{FTS/STS} 训练位置 \\
RSSI & \eng{rssiDbm} & 整符号功率；\modref{6.2.2.1}{波形和符号} 符号粒度 \\
RSRQ & \eng{rsrq} & \modref{6.5.5.1}{RSRP}/\modref{6.5.5.2}{RSSI} 同符号 \\
SINR & \eng{sinr} & 同上；分母$\le0$ 判错 \\
CBRA & \eng{cbra} & 门限 $-87/-62$\,dBm；DT/帧长 \\
PMS CSI & \eng{pmsIqLinear}/\eng{pmsRplDbm} & 模式1/2；\modref{6.2.3.9}{PMS波形} \\
时间差 & $T$ 或 $T_1,T_2$ & $t_1{\ldots}t_6$、角色；下游 \modref{6.6}{位置测量流程} \\
AoA & 水平/垂直角 & 阵列快照；坐标系约定 \\
\end{longtable}

%----------------------------------------------------------------------
\part{6.5.6}
%----------------------------------------------------------------------

\section{6.5.6 调制}
\textbf{API}：\eng{slbSelectModType} / \eng{slbGenerateModulatedSymbols}

\begin{longtable}{L{3.5cm}L{11.7cm}}
\caption{6.5.6}\label{tab:656} \\
\toprule \textbf{角色} & \textbf{内容} \\ \midrule
\endfirsthead
\toprule \textbf{角色} & \textbf{内容} \\ \midrule
\endhead
\bottomrule\endfoot
输出 &
\eng{modulatedSymbols[]}；
\eng{bitsPerSymbolM}；
\eng{symbolCount} \\
直接入参 &
L1 \eng{modType}；
比特流（$N_{\mathrm{bit}}\bmod m{=}0$） \\
他模块 &
数据面必须 \modref{6.5.7}{比特加扰}\eng{scrambledBits[]}；
控制面 \modref{6.2.6}{信道编码}/\modref{6.2.4}{物理层控制信息} \\
L0 & $1/\sqrt{2,10,42,170,682,2730}$；公式含整体负号 \\
下游 & \modref{6.2.3}{物理层信号} 映射；可选 \modref{6.5.8}{SC-FDMA}；\modref{6.5.4}{功率控制} 再缩放 \\
\end{longtable}

%----------------------------------------------------------------------
\part{6.5.7--6.5.9}
%----------------------------------------------------------------------

\section{6.5.7 比特加扰}
\textbf{API}：\eng{slbCalculateScrambleCinit} / \eng{slbScrambleBits}

\begin{longtable}{L{3.5cm}L{11.7cm}}
\caption{6.5.7}\label{tab:657} \\
\toprule \textbf{角色} & \textbf{内容} \\ \midrule
\endfirsthead
\toprule \textbf{角色} & \textbf{内容} \\ \midrule
\endhead
\bottomrule\endfoot
输出 & \eng{scrambleCinit}；\eng{scrambledBits[]} \\
直接入参 & 编码比特；业务类型；$n_f$、$N_{\mathrm{G-PID}}$ 等 \\
他模块 & \modref{6.5.2}{Gold序列}\eng{goldBits[]}；\modref{6.2.5}{物理层数据}/\modref{6.2.6}{信道编码} 比特 \\
隔离 & 与 6.5.3 / 6.5.9.1 的 $c_{\mathrm{init}}$ 禁止混流 \\
\end{longtable}

\section{6.5.8 SC-FDMA}
\textbf{API}：\eng{slbGenerateScFdmaPrecoded}

\begin{longtable}{L{3.5cm}L{11.7cm}}
\caption{6.5.8}\label{tab:658} \\
\toprule \textbf{角色} & \textbf{内容} \\ \midrule
\endfirsthead
\toprule \textbf{角色} & \textbf{内容} \\ \midrule
\endhead
\bottomrule\endfoot
输出 & \eng{scFdmaSymbols[]}；$M$；$L{=}M_{\mathrm{symb}}/M$ \\
直接入参 & $M_{\mathrm{SC}}^{\mathrm{Total}}$；\modref{6.5.6}{调制}\eng{modulatedSymbols[]} \\
约束 & $M_{\mathrm{symb}}$ 整除 $M$；强制 $1/\sqrt{M}$；不做 IFFT/CP \\
\end{longtable}

\section{6.5.9 / 6.5.9.1 ZC 与组跳}
\textbf{API}：\eng{slbGenerateZcSequence} / \eng{slbCalculateZcGroupHop}

\begin{longtable}{L{3.5cm}L{11.7cm}}
\caption{6.5.9}\label{tab:659} \\
\toprule \textbf{角色} & \textbf{内容} \\ \midrule
\endfirsthead
\toprule \textbf{角色} & \textbf{内容} \\ \midrule
\endhead
\bottomrule\endfoot
ZC 输出 & \eng{zcBaseSequence[]}；\eng{zcCyclicShifted[]} \\
ZC 入参 & $u,v,\alpha,m,M_{\mathrm{SC}}^{\mathrm{RS}}$；$m{=}2$ 必须查表 \\
组跳输出 & \eng{groupHopCinit}；\eng{groupIndexU} \\
组跳入参 & $n_{\mathrm{ID}}^X$、$n_S$；\modref{6.5.2}{Gold序列}\eng{goldBits[]} \\
\end{longtable}

\section{嵌套依赖示意}
\begin{center}\small
6.5.2 Gold $\rightarrow$ 6.5.3 / 6.5.7 / 6.5.9.1（三 $c_{\mathrm{init}}$ 隔离）\\[0.3em]
6.5.5.1 RSRP $\rightarrow$ 6.5.4 功控；同符号 RSSI $\rightarrow$ RSRQ/SINR\\[0.3em]
6.5.7 $\rightarrow$ 6.5.6 $\rightarrow$（可选 6.5.8）$\rightarrow$ 映射/功控
\end{center}

\end{document}
